\documentclass[main.tex]{subfiles}
\begin{document}

\section*{Traffic congestion in Cape Town and the alleviation
potential of DRT (\textcolor{red}{split up: 1) Situation in Cape Town, 2) Concept of Congestion and why congestion arises and how to measure it? 3) Drt as promising alternative or curse due to PT cannibalization? })}


Cape Town is the second-largest city in South Africa, and one of the most advanced cities in Africa. Its transportation system is characterised by high levels of household income spent on transport, long walking distances, waiting times and congestion \citep{StateofCT, jennings2015finding}. This is, at least in part, the result of urban sprawl, where most residents live on the periphery \citep{Turok2011, jennings2015finding} because of apartheid spatial planning \citep{Liphoto2025}. Transport therefore occurs in a tidal pattern with most public transport trips occurring to and from the city centre \citep{Turok2011, jennings2015finding}. Public transport waiting times, and travel times are generally high in the city. \citealp{hitge2015comparison} note the large difference in travel times between public transport and private transport in the city, affecting levels of accessibility. Compared to many other cities on the continent, public transport in South African cities is characterised by a unique modal mix of scheduled fixed-route public transport systems (passenger rail services and scheduled bus services) and a paratransit system, comprising unscheduled, semi-regulated modes \citep{Ryseck2024}. The paratransit modes cater to more than 70\,\% of the public transport modal split in the city despite service quality issues, such as overcrowding, safety and unreliability, amongst others \citep{Ferro2024}. In addition, ride-hailing services including Uber and Bolt, also serve a portion of the population. 

A transport survey in the Western Cape Province by Statistics South Africa reveals that many Cape Town residents (14.8\,\%) cite congestion as the main traffic problem  they experience \citep{StatsSAreport}. Travel time was the most common determinant of transport mode choice (32.5\,\%) in the 2013 National Household Travel Survey \citep{StatsSA14}. These community survey results provide evidence that the experiences of city residents point to problems of congestion and a lack of public transport options. Public transport in Cape Town faces many challenges, such as, “overcrowding, slow journey times, poor modal integration, a lack of off-peak services, manual ticketing and an irregular timetable” \citep{TDCIP}.

The problem of congestion is recognised by the transport authorities as costly in terms of time, money, and environmental impacts, and they have developed a “Congestion Management Plan” \citep{TDCIP}. \citealp{jennings2015finding} also notes the long travel times experienced in Cape Town, and the problem of congestion experienced as traffic converges from surrounding areas on the CBD. 2013 data showed that travel took 71\,\% longer than it would without congestion - these are large costs for a city to bear \citep{TDA}. The economic costs of congestion in Cape Town have been estimated to be equivalent to 2.8 billion Rand or USD 160\,m in 2018 \citep{IOL}. Infrastructure projects to deal with congestion alone have cost 192 m Rand or USD 11m, whilst  the total for capacity improvement and running transport in the city is 6.4 billion Rand or USD362\,m over 3 years \citep{TCT}. These spending items reveal the city's estimation of the cost of congestion that they can address with their budget and the importance of efficient transport to the city.
 
Moreover, accessibility has been raised as a key framework for transport design and justice \citep{Access}. It is especially important in Cape Town, which has a history of discriminatory exclusionary policy that impacted city design and thus transport corridors \citep{Liphoto2025}. Unemployment and spatial segregation in the city form an exclusionary combination that needs to be addressed with better transport and land use planning. \citealp{Moyo} found that middle-income areas have more mixed land use, while poor areas remain homogenous. 

The city of Cape Town is aiming to increase road capacity and public transport particularly for low-income areas in the coming years \citep{IDP}. The problem of congestion will likely persist even with these interventions, and DRT could be a useful complement to these efforts, especially as they do not focus on affluent areas, which contribute large volumes of vehicles. \citealp{byala2021redesigning} note that new mobility such as DRT can be incorporated into redesigning bus systems, this could be a beneficial exercise in the future of Cape Town’s traffic design as the city changes and grows and demand for efficient transport increases.

DRT systems could be a helpful addition to bridge the time and cost gap between public and private transport options in Cape Town, a space where shared mobility options are suitable. The basic idea of DRT is that of transport provided flexibly in response to user demand signals. \citealp{butler2020smart} defines DRT as being able to change routes and or schedule depending on demand. DRT has been classified into several forms, and it overlaps with the categories of paratransit and shared mobility \citep{butler2020smart}. The key to sharing is a higher utilization of goods by replacing permanent individual ownership with temporary on-demand access \citep{botsman2010s}. Whether it is classified together with sharing or not, the main idea is that DRT vehicles service many more users than a private car would \citep{munzel2018carsharing, butler2020smart}. It can either be area-based or line-based, with door-to-door stops \citep{vansteenwegen2022survey} or set stop points when it is stop-based. It can be ad-hoc, similar to existing minibus taxis, or centrally coordinated, such as the MOIA service in Hamburg, Germany \citep{KagerbauerKostorzWilkes2021_1000141282}. In this paper, DRT refers to some real-time centrally coordinated DRT using app-based technology, as such technology is well suited for dynamic requests and routing \citep{zwick2021ride}. As DRT is known to specifically work in mixed land-use scenarios, it would make sense to first investigate the more affluent areas of Cape Town for DRT viability \citep{yu2019exploring}. Design of DRT could be incorporated or complement future transit system redesign in Cape Town \citep{byala2021redesigning}.

The general need to improve transport in cities has been termed ‘smart urban mobility’\citep{butler2020smart} and is not limited to Cape Town. A literature review on urban mobility scenarios by \citealp{miskolczi2021urban} found that the most likely outcome is a slow shift to emerging technologies of ride sharing, electric and self-driving vehicles by 2030. In this context, travel time is a key factor. Travel time is essential in decisions to use transport modes such as DRT \citep{reck2020shared}, and high fares do not necessarily deter use of public transport \citep{tyrinopoulos2013factors}. Thus, prioritising efficient transport for lower congestion and time costs is a good priority, and aligns well with the expectation of the literature of a shift to an increased use of DRT and similar emerging technologies.

However, the real-world implications of a large-scale DRT service on transportation system determinants like congestion remain unclear. \citealp{butler2020smart} note the effects of DRT on congestion can be both positive and negative, because congestion is a complex phenomenon \citep{arnott1994economics}. Congestion frustrates drivers with high fuel costs and travel times. Those who use public transport often have long queues, overcrowding, and delays \citep{jennings2015finding}. DRT can decrease cars per passenger and increase traffic flow because of adaptive planning, but pick up and drop-off zones may also increase congestion, and vehicle numbers can increase with higher accessibility or preference for low-occupancy vehicles. These issues should be monitored in pilot studies after the simulation and tests are completed. Furthermore, land use and urban planning must be designed to complement purely mobility-focused solutions to obtain the most efficiency \citep{jabbarpour2018applications}. Small decreases in travel demand can have significant impacts on alleviating congestion and improving travel time \citep{vickrey1969congestion}. This, however, requires particular conditions, as congestion often occurs not only because of physical bottlenecks \citep{arnott1994economics}. 

In that vein, a report by the International Network for Transport and Accessibility in Low Income Communities (INTALInC) highlights the need for comparative studies that match local contexts \citep{lucas2019transport}. Simulations with locally calibrated data fill this gap and harness the power of big data to help understand and solve social problems. This would also be in line with \citealp{vanderschuren2018ride}, who evaluated the potential market for ridesharing in Cape Town using the 2013 national transport survey data. Their research, however, is constrained by a lack of quantitative data and relies on expert opinion alone. Using ABM to simulate the city transport system, however, addresses this challenge of limited data availability and opens up the possibility for an evaluation of experimental policies such as the introduction of large-scale DRT.  

% and we can investigate the author's conclusions that low income households would switch to ridesharing solutions and DRT could have a vital role to play in this regard , hence serving transport demand in the city.

A  simulation of Hamburg, Germany's MOIA service, a ridepooling service for up to five passengers in fully electric autonomous vehicles and varying fleet sizes and service areas \citep{KagerbauerKostorzWilkes2021_1000141282} finds that increasing the use of DRT has a positive effect on the use of public transport rather than 'cannibalising' (\textcolor{red}{this is the first time cannibalization is mentioned, should we explain the theoretical t Framework of DRT-Transit Integration vs. Cannibalization?} it, which is often a hypothetical argument against introducing DRT. Furthermore in their research the large fleets with autonomous vehicles lead to the lowest costs \citep{KagerbauerKostorzWilkes2021_1000141282}. Others \cite{Basu.2018} find that ride-pooling services serve as a complement to public transport and increase the efficiency of the overall system by serving as a feeder to the transit system. On the contrary, it has been highlighted that introducing on-demand services comes with the potential threat of public transport cannibalized by shared DRT \cite{Oh.2020}. Simulations for Jerusalem \cite{BenDor.2022}, Berlin \cite{Kaddoura.2020} and the Rouen Normandie metropolitan area \cite{Vosooghi.2019} also indicate that cannibalization of the existing public transport sector by DRT must be avoided. Therefore, the introduction of 'push factors'—measures that make private car use less attractive (e.g., higher parking fees or congestion tolls)—might be necessary. As presented in this literature overview, there are many studies on DRT simulations and outcomes, but none focus on congestion and in particular not in the context of the Global South. This study attempts to fill this research gap.

%A research theme losely related to DRT is Mobility-as-a-Service (MaaS). The idea behind MaaS is to have integrated transport offerings within city or regional systems \citealp{Loubser} Such motivation-based policies offer higher efficiency and lower emissions \citep{maggi2021price}. DRT is often seen as playing a key role in MaaS. Since the simulation in this work will primarily have time and price-based switching, it can be interpreted as a lower bound of modal switching without preference influencing policy or marketing being implemented to increase the switching.




\end{document}